\chapter{Dihedral Groups}

\section{Groups Generated by Two Involutions}

\begin{definition}[Involution]
An element of order $2$ in a group is called an \textbf{involution}.
\end{definition}

\begin{definition}[Dihedral group]
Let
\[
    G=\gen{a,b},
    \qquad
    |a|=|b|=2.
\]
Put
\[
    x=ab.
\]
If $|x|=n<\infty$, then $G$ is called a \textbf{dihedral group} of order $2n$, and we write
\[
    G=D_{2n}.
\]
If $|x|=\infty$, then $G$ is called the \textbf{infinite dihedral group}.
\end{definition}

\begin{example}
Let
\[
    a=
    \begin{pmatrix}
        0 & 1\\
        1 & 0
    \end{pmatrix},
    \qquad
    b=
    \begin{pmatrix}
        0 & 2\\
        \frac12 & 0
    \end{pmatrix}
    \in \operatorname{GL}_2(\mathbb R).
\]
Then
\[
    a^2=b^2=I_2,
\]
so both $a$ and $b$ are involutions. Moreover,
\[
    ab=
    \begin{pmatrix}
        \frac12 & 0\\
        0 & 2
    \end{pmatrix}.
\]
Hence
\[
    (ab)^m=
    \begin{pmatrix}
        2^{-m} & 0\\
        0 & 2^m
    \end{pmatrix}\neq I_2
    \qquad(m\neq 0).
\]
Thus $|ab|=\infty$, and therefore
\[
    G=\gen{a,b}
\]
is an infinite dihedral group.
\end{example}

\begin{proposition}[Normal form]
Let
\[
    G=\gen{a,b},
    \qquad
    |a|=|b|=2,
\]
and put $x=ab$. Then every element of $G$ has the form
\[
    x^i
    \qquad\text{or}\qquad
    x^i b
\]
for some $i\in\mathbb Z$.
\end{proposition}

\begin{proof}
Since
\[
    a^2=b^2=1,
\]
any word in $a$ and $b$ can be reduced by deleting adjacent subwords $aa$ and $bb$.

Hence every reduced word is alternating. Thus the possible words are of the following types:
\[
    a,\ ab,\ aba,\ abab,\ldots
\]
and
\[
    b,\ ba,\ bab,\ baba,\ldots .
\]
Equivalently,
\[
    G=
    \left\{
        (ab)^i,\ (ba)^j,\ (ab)^i a,\ (ba)^j b
        : i,j\geq 0
    \right\}.
\]

Now
\[
    (ba)^j
    =
    b(ab)^j b
    =
    \bigl((ab)^j\bigr)^{-1}
    =
    x^{-j}.
\]
Also,
\[
    (ab)^i a
    =
    (ab)^i a b b
    =
    (ab)^{i+1}b
    =
    x^{i+1}b,
\]
and
\[
    (ba)^j b
    =
    b(ab)^j
    =
    x^{-j}b.
\]
Therefore every element of $G$ is of the form
\[
    x^i
    \qquad\text{or}\qquad
    x^i b
\]
for some $i\in\mathbb Z$.
\end{proof}

\begin{theorem}
Let
\[
    G=\gen{a,b},
    \qquad
    |a|=|b|=2.
\]
Then $G$ is finite if and only if $ab$ has finite order.

Moreover, if
\[
    |ab|=n,
\]
then
\[
    |G|=2n.
\]
\end{theorem}

\begin{proof}
Put
\[
    x=ab.
\]
If $G$ is finite, then every element of $G$ has finite order, so $x=ab$ has finite order.

Conversely, suppose that $|x|=n<\infty$. By the normal form above,
\[
    G=\{x^i,\ x^i b: i\in\mathbb Z\}.
\]
Since $x^n=1$, we may reduce the exponent $i$ modulo $n$. Hence
\[
    G=
    \{1,x,x^2,\ldots,x^{n-1}\}
    \cup
    \{b,xb,x^2b,\ldots,x^{n-1}b\}.
\]
Thus $G$ is finite and $|G|\leq 2n$.

Let
\[
    H=\gen{x}.
\]
Then $|H|=n$. We claim that
\[
    H\cap Hb=\varnothing.
\]
Indeed, if $x^i=x^j b$, then
\[
    b=x^{j-i}\in H.
\]
It follows that both $a=xb$ and $b$ lie in the cyclic group $H$. Since $a$ and $b$ are both involutions, this would force $a=b$, and hence
\[
    x=ab=1,
\]
contradicting $|x|=n$ unless $n=1$. In the case $n=1$, we have $a=b$ and $H=\{1\}$, so again $b\notin H$. Therefore $H$ and $Hb$ are disjoint.

Hence
\[
    |G|=|H|+|Hb|=n+n=2n.
\]
\end{proof}

\section{Semidirect Product Description}

\begin{definition}[Internal semidirect product]
Let $G$ be a group, and let $A,B\leq G$. We say that $G$ is an \textbf{internal semidirect product} of $A$ by $B$ if
\begin{enumerate}[label=(\arabic*)]
    \item $A\lhd G$;
    \item $A\cap B=1$;
    \item $G=AB$.
\end{enumerate}
In this case, we write
\[
    G=A:B
    \qquad\text{or}\qquad
    G=A\rtimes B.
\]
\end{definition}

\begin{proposition}
Let
\[
    G=D_{2n}=\gen{a,b},
    \qquad
    |a|=|b|=2,
    \qquad
    |ab|=n.
\]
Put
\[
    x=ab.
\]
Then
\[
    G=\gen{x}:\gen{b}.
\]
In particular,
\[
    D_{2n}\cong C_n:C_2.
\]
\end{proposition}

\begin{proof}
We have
\[
    x=ab,
    \qquad
    a=xb.
\]
Therefore
\[
    G=\gen{a,b}=\gen{x,b}.
\]

Moreover,
\[
    bxb
    =
    b(ab)b
    =
    ba
    =
    (ab)^{-1}
    =
    x^{-1}.
\]
Hence $b$ normalizes $\gen{x}$, and so
\[
    \gen{x}\lhd G.
\]
Also,
\[
    \gen{x}\cap\gen{b}=1,
\]
and
\[
    G=\gen{x}\gen{b}.
\]
Thus
\[
    G=\gen{x}:\gen{b}.
\]
Since
\[
    \gen{x}\cong C_n,
    \qquad
    \gen{b}\cong C_2,
\]
we get
\[
    D_{2n}\cong C_n:C_2.
\]
Here the element $b$ acts on $C_n=\gen{x}$ by inversion:
\[
    x\longmapsto x^{-1}.
\]
\end{proof}

\section{Automorphisms of \texorpdfstring{$D_{2n}$}{D2n} for Odd \texorpdfstring{$n$}{n}}

Throughout this section, write
\[
    D_{2n}
    =
    \gen{g,b}
    =
    \gen{g}:\gen{b},
\]
where
\[
    |g|=n,
    \qquad
    |b|=2,
    \qquad
    bgb=g^{-1}.
\]
Thus
\[
    D_{2n}
    =
    \{g^i,\ g^i b:0\leq i\leq n-1\}.
\]

\begin{theorem}[Automorphism group of $D_{2n}$ for $n$ odd]
Assume that $n$ is odd. Then every automorphism of $D_{2n}$ has the form
\[
    \theta_{i,j}:
    \begin{cases}
        g\longmapsto g^i,\\
        b\longmapsto g^j b,
    \end{cases}
\]
where
\[
    \gcd(i,n)=1,
    \qquad
    j\in\mathbb Z/n\mathbb Z.
\]
Consequently,
\[
    |\Aut(D_{2n})|=n\varphi(n),
\]
and
\[
    \Aut(D_{2n})
    \cong
    C_n:\Aut(C_n).
\]
\end{theorem}

\begin{proof}
Let
\[
    \theta\in\Aut(D_{2n}).
\]
Since $|g|=n$ and $n$ is odd, the element $\theta(g)$ also has odd order $n$.

But every element outside $\gen{g}$ has the form $g^j b$, and
\[
    (g^j b)^2
    =
    g^j b g^j b
    =
    g^j g^{-j}b^2
    =
    1.
\]
Hence every element outside $\gen{g}$ has order $2$. Therefore
\[
    \theta(g)\in\gen{g}.
\]
So
\[
    \theta(g)=g^i
\]
for some $i$ with
\[
    \gcd(i,n)=1.
\]

Next, $\theta(b)$ has order $2$. Since $n$ is odd, the only elements of order $2$ in $D_{2n}$ are the elements
\[
    g^j b,
    \qquad
    0\leq j\leq n-1.
\]
Therefore
\[
    \theta(b)=g^j b
\]
for some $j\in\mathbb Z/n\mathbb Z$.

Conversely, for any $i,j$ satisfying
\[
    \gcd(i,n)=1,
    \qquad
    j\in\mathbb Z/n\mathbb Z,
\]
define
\[
    \theta_{i,j}(g)=g^i,
    \qquad
    \theta_{i,j}(b)=g^j b.
\]
We check that the defining relations are preserved:
\[
    (g^i)^n=1,
\]
and
\[
    (g^j b)^2
    =
    g^j b g^j b
    =
    g^j g^{-j}b^2
    =
    1.
\]
Also,
\[
    (g^j b)g^i(g^j b)^{-1}
    =
    g^j b g^i b g^{-j}
    =
    g^j g^{-i}g^{-j}
    =
    g^{-i}.
\]
Thus $\theta_{i,j}$ gives a homomorphism $D_{2n}\to D_{2n}$.

Since $\gcd(i,n)=1$, the element $g^i$ generates $\gen{g}$. Hence the image of $\theta_{i,j}$ contains $\gen{g}$. It also contains $g^j b$, and therefore contains
\[
    b=g^{-j}(g^j b).
\]
Thus $\theta_{i,j}$ is surjective, hence an automorphism.

Therefore
\[
    \Aut(D_{2n})
    =
    \{\theta_{i,j}:\gcd(i,n)=1,\ j\in\mathbb Z/n\mathbb Z\},
\]
so
\[
    |\Aut(D_{2n})|
    =
    n\varphi(n).
\]

Finally, define
\[
    N=\{\theta_{1,j}:j\in\mathbb Z/n\mathbb Z\}.
\]
Then
\[
    N\cong C_n.
\]
Also define
\[
    U=\{\theta_{i,0}:\gcd(i,n)=1\}.
\]
Then
\[
    U\cong\Aut(C_n).
\]
Moreover,
\[
    \theta_{i,j}\circ\theta_{k,\ell}
    =
    \theta_{ik,\ j+i\ell}.
\]
This shows that
\[
    \Aut(D_{2n})
    =
    N:U
    \cong
    C_n:\Aut(C_n).
\]
\end{proof}

\begin{corollary}
Let $p$ be an odd prime. Then
\[
    \Aut(D_{2p})
    \cong
    C_p:C_{p-1}.
\]
More generally, for $e\geq 1$,
\[
    \Aut(D_{2p^e})
    \cong
    C_{p^e}:C_{p^{e-1}(p-1)}.
\]
\end{corollary}

\begin{proof}
By the theorem,
\[
    \Aut(D_{2n})
    \cong
    C_n:\Aut(C_n)
\]
for odd $n$.

If $n=p$, then
\[
    \Aut(C_p)\cong C_{p-1}.
\]
Hence
\[
    \Aut(D_{2p})
    \cong
    C_p:C_{p-1}.
\]

If $n=p^e$, then
\[
    \Aut(C_{p^e})
    \cong
    C_{p^{e-1}(p-1)}.
\]
Therefore
\[
    \Aut(D_{2p^e})
    \cong
    C_{p^e}:C_{p^{e-1}(p-1)}.
\]
\end{proof}

\begin{theorem}[Automorphism group of $D_{2n}$ for $n$ even]
    Assume that $n$ is even and $n\geq 4$. Then every automorphism of $D_{2n}$ has the form
    \[
        \theta_{i,j}:
        \begin{cases}
            g\longmapsto g^i,\\
            b\longmapsto g^j b,
        \end{cases}
    \]
    where
    \[
        \gcd(i,n)=1,
        \qquad
        j\in\mathbb Z/n\mathbb Z.
    \]
    Consequently,
    \[
        |\Aut(D_{2n})|=n\varphi(n),
    \]
    and
    \[
        \Aut(D_{2n})
        \cong
        C_n:\Aut(C_n).
    \]
    \end{theorem}
    
    \begin{proof}
    Let
    \[
        D_{2n}=\gen{g,b},
        \qquad
        |g|=n,
        \qquad
        |b|=2,
        \qquad
        bgb=g^{-1}.
    \]
    Then
    \[
        D_{2n}
        =
        \{g^i,\ g^i b:0\leq i\leq n-1\}.
    \]
    
    Let
    \[
        \theta\in\Aut(D_{2n}).
    \]
    Since $g$ has order $n$, the element $\theta(g)$ also has order $n$.
    
    But every element outside $\gen{g}$ has the form $g^j b$, and
    \[
        (g^j b)^2
        =
        g^j b g^j b
        =
        g^j g^{-j}b^2
        =
        1.
    \]
    Hence every element outside $\gen{g}$ has order $2$.
    
    Since $n\geq 4$, an element of order $n$ cannot lie outside $\gen{g}$. Therefore
    \[
        \theta(g)\in\gen{g}.
    \]
    So
    \[
        \theta(g)=g^i
    \]
    for some $i$ satisfying
    \[
        \gcd(i,n)=1.
    \]
    
    Next, $\theta(b)$ has order $2$. Since $\theta$ is surjective, $\theta(b)$ cannot lie in $\gen{g}$; otherwise the image of $\theta$ would be contained in $\gen{g}$. Hence
    \[
        \theta(b)=g^j b
    \]
    for some
    \[
        j\in\mathbb Z/n\mathbb Z.
    \]
    
    Conversely, take $i,j$ such that
    \[
        \gcd(i,n)=1,
        \qquad
        j\in\mathbb Z/n\mathbb Z.
    \]
    Define
    \[
        \theta_{i,j}(g)=g^i,
        \qquad
        \theta_{i,j}(b)=g^j b.
    \]
    We check that the defining relations are preserved.
    
    First,
    \[
        (g^i)^n=1.
    \]
    Second,
    \[
        (g^j b)^2
        =
        g^j b g^j b
        =
        g^j g^{-j}b^2
        =
        1.
    \]
    Finally,
    \[
        (g^j b)g^i(g^j b)^{-1}
        =
        g^j b g^i b g^{-j}
        =
        g^j g^{-i}g^{-j}
        =
        g^{-i}.
    \]
    Thus $\theta_{i,j}$ defines a homomorphism
    \[
        D_{2n}\longrightarrow D_{2n}.
    \]
    
    Since $\gcd(i,n)=1$, the element $g^i$ generates $\gen{g}$. Hence the image of $\theta_{i,j}$ contains $\gen{g}$. Also, it contains $g^j b$, so it contains
    \[
        b=g^{-j}(g^j b).
    \]
    Therefore the image contains both $g$ and $b$, and hence $\theta_{i,j}$ is surjective. Since $D_{2n}$ is finite, $\theta_{i,j}$ is an automorphism.
    
    Thus
    \[
        \Aut(D_{2n})
        =
        \{\theta_{i,j}:\gcd(i,n)=1,\ j\in\mathbb Z/n\mathbb Z\}.
    \]
    There are $\varphi(n)$ choices for $i$ and $n$ choices for $j$, so
    \[
        |\Aut(D_{2n})|=n\varphi(n).
    \]
    
    Moreover, put
    \[
        N=\{\theta_{1,j}:j\in\mathbb Z/n\mathbb Z\},
        \qquad
        U=\{\theta_{i,0}:\gcd(i,n)=1\}.
    \]
    Then
    \[
        N\cong C_n,
        \qquad
        U\cong\Aut(C_n).
    \]
    Also,
    \[
        \theta_{i,j}\circ\theta_{k,\ell}
        =
        \theta_{ik,\ j+i\ell}.
    \]
    Hence
    \[
        \Aut(D_{2n})
        =
        N:U
        \cong
        C_n:\Aut(C_n).
    \]
    \end{proof}
    
    \begin{corollary}
    Let $m\geq 2$. Then
    \[
        \Aut(D_{4m})
        \cong
        C_{2m}:\Aut(C_{2m}),
    \]
    and hence
    \[
        |\Aut(D_{4m})|
        =
        2m\,\varphi(2m).
    \]
    \end{corollary}
    
    \begin{proof}
    This follows from the theorem by taking
    \[
        n=2m.
    \]
    \end{proof}
    
    \begin{corollary}
    Let $e\geq 2$. Then
    \[
        \Aut(D_{2^{e+1}})
        \cong
        C_{2^e}:\Aut(C_{2^e}).
    \]
    In particular,
    \[
        |\Aut(D_{2^{e+1}})|
        =
        2^e\varphi(2^e)
        =
        2^{2e-1}.
    \]
    More explicitly,
    \[
        \Aut(C_{2^e})
        \cong
        \begin{cases}
            C_2, & e=2,\\
            C_2\times C_{2^{e-2}}, & e\geq 3.
        \end{cases}
    \]
    Therefore
    \[
        \Aut(D_8)\cong C_4:C_2,
    \]
    and for $e\geq 3$,
    \[
        \Aut(D_{2^{e+1}})
        \cong
        C_{2^e}:(C_2\times C_{2^{e-2}}).
    \]
    \end{corollary}
    
    \begin{remark}
    The case $e=1$ is exceptional. In that case,
    \[
        D_4\cong C_2\times C_2,
    \]
    so
    \[
        \Aut(D_4)\cong S_3.
    \]
    Thus the formula above starts from $e\geq 2$.
    \end{remark}

\section{Conjugacy Classes in \texorpdfstring{$D_{2n}$}{D2n}}

Let
\[
    D_{2n}
    =
    \gen{g}:\gen{b}
    =
    \{g^i,\ g^i b:0\leq i\leq n-1\},
\]
where
\[
    |g|=n,
    \qquad
    |b|=2,
    \qquad
    bgb=g^{-1}.
\]

\begin{lemma}
For every integer $i$,
\[
    bg^i b=g^{-i}.
\]
Consequently,
\[
    b^{-1}g^i b=g^{-i}.
\]
Also, for all integers $i,j$,
\[
    g^j(g^i b)g^{-j}
    =
    g^{i+2j}b.
\]
\end{lemma}

\begin{proof}
Since
\[
    bgb=g^{-1},
\]
we get
\[
    bg^i b=(bgb)^i=(g^{-1})^i=g^{-i}.
\]
Because $b^{-1}=b$, this also gives
\[
    b^{-1}g^i b=g^{-i}.
\]

For the second identity,
\[
    g^j(g^i b)g^{-j}
    =
    g^{i+j} b g^{-j}.
\]
Using
\[
    bg^{-j}=g^j b,
\]
we get
\[
    g^j(g^i b)g^{-j}
    =
    g^{i+j}g^j b
    =
    g^{i+2j}b.
\]
\end{proof}

\subsection{The case where \texorpdfstring{$n$}{n} is odd}

\begin{theorem}[Conjugacy classes for odd $n$]
Assume that $n$ is odd. Then the conjugacy classes of $D_{2n}$ are
\[
    \{1\},
\]
\[
    \{g^i,g^{-i}\},
    \qquad
    1\leq i\leq \frac{n-1}{2},
\]
and
\[
    \{b,gb,g^2b,\ldots,g^{n-1}b\}.
\]
Equivalently, the elements of odd order have conjugacy classes
\[
    \{g^i,g^{-i}\},
\]
and all elements of even order are conjugate.
\end{theorem}

\begin{proof}
First consider the powers of $g$. Since $\gen{g}$ is abelian, conjugation by powers of $g$ fixes every $g^i$. Conjugation by $b$ gives
\[
    b^{-1}g^i b
    =
    g^{-i}.
\]
Therefore the conjugacy class of $g^i$ is
\[
    \{g^i,g^{-i}\}.
\]
Since $n$ is odd,
\[
    g^i=g^{-i}
    \iff
    g^{2i}=1
    \iff
    n\mid 2i
    \iff
    n\mid i.
\]
Thus the identity is its own conjugacy class, and the remaining rotations occur in pairs
\[
    \{g^i,g^{-i}\}.
\]

Now consider the elements outside $\gen{g}$. They are
\[
    g^i b,
    \qquad
    0\leq i\leq n-1.
\]
Each has order $2$, because
\[
    (g^i b)^2
    =
    g^i b g^i b
    =
    g^i g^{-i}b^2
    =
    1.
\]
So $D_{2n}$ has exactly $n$ elements of order $2$, namely
\[
    g^i b,
    \qquad
    0\leq i\leq n-1.
\]

We claim that all these elements are conjugate. Indeed,
\[
    g^j b g^{-j}
    =
    g^{2j}b.
\]
Since $n$ is odd, the map
\[
    j\longmapsto 2j
    \pmod n
\]
is a bijection of $\mathbb Z/n\mathbb Z$. Hence the elements
\[
    g^j b g^{-j},
    \qquad
    0\leq j\leq n-1,
\]
are exactly
\[
    b,gb,g^2b,\ldots,g^{n-1}b.
\]
Thus all elements of the form $g^i b$ are conjugate.

Therefore, when $n$ is odd, the elements of odd order are precisely the rotations, whose conjugacy classes are
\[
    \{g^i,g^{-i}\},
\]
and the elements of even order are precisely the reflections $g^i b$, which form a single conjugacy class.
\end{proof}

\begin{remark}
Equivalently, since $2$ is invertible modulo $n$, choose $t$ such that
\[
    2t\equiv 1\pmod n.
\]
Then
\[
    g^t(g^i b)g^{-t}
    =
    g^{i+2t}b
    =
    g^{i+1}b.
\]
Hence consecutive reflections
\[
    g^i b
    \qquad\text{and}\qquad
    g^{i+1}b
\]
are conjugate, and therefore all reflections are conjugate.
\end{remark}

\subsection{The case where \texorpdfstring{$n$}{n} is even}

\begin{theorem}[Conjugacy classes for even $n$]
    Assume that $n$ is even. Then the conjugacy classes of $D_{2n}$ are
    \[
        \{1\},
    \]
    \[
        \{g^{n/2}\},
    \]
    \[
        \{g^i,g^{-i}\},
        \qquad
        1\leq i\leq \frac n2-1,
    \]
    and the following two conjugacy classes of reflections:
    \[
        \{b,g^2b,g^4b,\ldots,g^{n-2}b\},
    \]
    \[
        \{gb,g^3b,g^5b,\ldots,g^{n-1}b\}.
    \]
    Equivalently, the reflections split into two conjugacy classes according to whether the exponent of $g$ is even or odd.
    \end{theorem}
    
    \begin{proof}
    Recall that
    \[
        D_{2n}
        =
        \gen{g}:\gen{b}
        =
        \{g^i,\ g^i b:0\leq i\leq n-1\},
    \]
    where
    \[
        |g|=n,
        \qquad
        |b|=2,
        \qquad
        bgb=g^{-1}.
    \]
    
    First consider the powers of $g$. Since $\gen{g}$ is abelian, conjugation by powers of $g$ fixes every $g^i$. On the other hand,
    \[
        b^{-1}g^i b
        =
        bg^i b
        =
        g^{-i}.
    \]
    Therefore the conjugacy class of $g^i$ is contained in
    \[
        \{g^i,g^{-i}\}.
    \]
    Moreover, $b$ conjugates $g^i$ to $g^{-i}$, so the conjugacy class of $g^i$ is exactly
    \[
        \{g^i,g^{-i}\},
    \]
    unless
    \[
        g^i=g^{-i}.
    \]
    
    Now
    \[
        g^i=g^{-i}
        \iff
        g^{2i}=1
        \iff
        n\mid 2i.
    \]
    Since $n$ is even, this happens precisely when
    \[
        i\equiv 0\pmod n
        \qquad\text{or}\qquad
        i\equiv \frac n2\pmod n.
    \]
    Hence
    \[
        \{1\}
    \]
    and
    \[
        \{g^{n/2}\}
    \]
    are single conjugacy classes. The remaining rotations occur in pairs
    \[
        \{g^i,g^{-i}\},
        \qquad
        1\leq i\leq \frac n2-1.
    \]
    
    Now consider the elements outside $\gen{g}$, namely the reflections
    \[
        g^i b,
        \qquad
        0\leq i\leq n-1.
    \]
    For any integer $j$, we have
    \[
        g^j(g^i b)g^{-j}
        =
        g^{i+2j}b.
    \]
    Thus conjugation by a power of $g$ changes the exponent $i$ by an even integer.
    
    Also,
    \[
        b^{-1}(g^i b)b
        =
        b(g^i b)b
        =
        bg^i
        =
        g^{-i}b.
    \]
    This again preserves the parity of $i$.
    
    Hence conjugation preserves the parity of the exponent $i$ in $g^i b$. Therefore a reflection with even exponent cannot be conjugate to a reflection with odd exponent.
    
    Conversely, suppose $i$ and $k$ have the same parity. Then
    \[
        k-i
    \]
    is even. Since $n$ is even, the congruence
    \[
        2j\equiv k-i\pmod n
    \]
    has a solution $j$. For such a $j$,
    \[
        g^j(g^i b)g^{-j}
        =
        g^{i+2j}b
        =
        g^k b.
    \]
    Therefore any two reflections whose exponents have the same parity are conjugate.
    
    Thus the reflections split into exactly two conjugacy classes:
    \[
        \{g^{2j}b:0\leq j\leq \frac n2-1\}
        =
        \{b,g^2b,g^4b,\ldots,g^{n-2}b\},
    \]
    and
    \[
        \{g^{2j+1}b:0\leq j\leq \frac n2-1\}
        =
        \{gb,g^3b,g^5b,\ldots,g^{n-1}b\}.
    \]
    This proves the theorem.
    \end{proof}
    
    \begin{corollary}
    Assume that $n$ is even. Then the involutions in $D_{2n}$ are
    \[
        g^{n/2}
    \]
    and
    \[
        g^i b,
        \qquad
        0\leq i\leq n-1.
    \]
    They split into three conjugacy classes:
    \[
        \{g^{n/2}\},
    \]
    \[
        \{b,g^2b,g^4b,\ldots,g^{n-2}b\},
    \]
    and
    \[
        \{gb,g^3b,g^5b,\ldots,g^{n-1}b\}.
    \]
    \end{corollary}
    
    \begin{proof}
    The element $g^{n/2}$ has order $2$, and every reflection $g^i b$ has order $2$ because
    \[
        (g^i b)^2
        =
        g^i b g^i b
        =
        g^i g^{-i}b^2
        =
        1.
    \]
    The conjugacy classes are exactly those described in the theorem.
    \end{proof}

